\section{n2. Exact smooth, quotient, coefficient-field and Frobenius-root current413 results independently reconstructed at their stated scopes.}
\textit{Supports: G331, n1.}

Independent reconstruction of current413's smooth, quotient, coefficient-field and root arguments. Supports are qualified G331 and own n1; current413 is DATA. The published BFJ application is ONLY a smooth affine CLOSED germ over an algebraically closed characteristic-zero field, dimension>=2. DVR and empty cases are n1. The complete immediate setting, modification/retraction chain1.10--1.15,2.1--2.8,3.1--3.14 and3.15--3.25, including the entire3.18 proof, were freshly read. Hironaka smooth characteristic-zero principalization, smooth multiplier subadditivity, thinness and smooth Izumi are qualified deeper foundations, not characteristic-free inputs.

For a model obstacle F=F composed with r\_pi0 let M=max(0,-min F). Nef divisors below its coefficients form a max-directed family after adjoining the barrier -M. Their coefficient supremum W lies in the weak closed nef cone and -M<=W<=0. BFJ3.11 supplies substantive continuity on every CLOSED simplex, and the full function H(v)=inf\_pi H(r\_pi v). Divisorial density and refinements of pi0 give H<=F on every support. G331's exact BFJ comparison gives every D competitor below H, so H is the raw envelope. Canonical B\_m=L-infinity(mW)=L-infinity(mZ(F)) contains a primary power. To identify it with n1's all-support A\_m, test (g)+mhat\textasciicircum{}N with N>=mM: its divisorial inequalities extend by closed-simplex continuity and cofinal retractions, and v(g)>=v((g)+mhat\textasciicircum{}N). Completion and exact primary contraction in G331 identify B\_m with A\_m in the local ring.

If W=0, F>=0 by density and the envelope is zero. Otherwise BFJ3.14 makes H strictly negative. Thinness is bounded only on the fixed compact complex; choose C>0 with A<=-CW there. The complete3.18 proof yields W\_t<=W<=(1-C/t)W\_(t-C) for real t>C: part(i) applies to real scalar multiples, and part(ii) is the same defining fixed-complex inequality for real t. Put t=m+C to obtain (1+C/m)W<=W\_m<=W. Closed simplices and all retraction infima give 0<=H-h\_m<=CM/m on EVERY support. Thus H belongs to continuous geometric C. Signed model functions are uniformly dense by the lattice ideal identities and BFJ1.14--1.15. N1's Lipschitz estimate gives CE for every continuous signed obstacle at exactly these hypotheses. No FG is used.

For a finite faithful group G acting locally on such a smooth S with finite excellent R=S\textasciicircum{}G, R is normal and Frac(S)\textasciicircum{}G=Frac(R), by orbit denominator products. The fixed-field degree proof is elementary: distinct automorphisms are linearly independent as functions, by subtracting a shortest relation evaluated at ax from one evaluated at x; an invertible evaluation matrix gives a basis over the fixed field and degree |G|. Orbit polynomials prove Galoisness. G331's normal-target coefficient trace has full fibre support and is the orbit measure sum\_sigma delta\_(w composed with sigma), even at nonzero supports. Hence every restriction fibre is a full G-orbit. The source greatest minorant for a pulled obstacle is invariant by uniqueness, thus positively descends through EVERY fibre and is the target greatest minorant. If B\_m is its stable canonical ideal, the coefficients J\_m of product\_sigma(sum\_i T\_i sigma(s\_i)), for finite R-module generators s\_i, lie in normal R and satisfy v(J\_m)=|G| w(B\_m) on every support. They are actual unit-or-primary ideals by G331. Their profiles descend the source approximants exactly.

For fixed signed primary/model obstacles with source error E/m, J\_m subset A\_(dm)\textasciicircum{}R, d=|G|. For k>=2d, let m=floor(k/d), r=dm. Grading with A\_1\textasciicircum{}(k-r) gives 0<=p\_R-h\_k\textasciicircum{}R<=E/floor(k/d)+dM\_1/k<=(2dE+dM\_1)/k. This is the precise conditional residue estimate, not same-degree contraction or a general finite-map envelope formula.

The separable coefficient extension argument is valid for a geometrically integral, geometrically normal positive-dimensional normal local domain R essentially of finite type over k with residue field EXACTLY k. For finite Galois E/k, R\_E=R tensor E is finite free and local since its closed fibre is the field E; the union R\_s over k\textasciicircum{}s is the corresponding Noetherian excellent normal local ring. Compatible semivaluations on its finite stages define all semivaluations on the union, so X\_s=lim\_E X\_E and restriction is compact surjective. At each finite stage G331's full-fibre trace and the fixed-field proof above give full Galois orbits. For two points in an infinite fibre, the sets of gamma in Gal(k\textasciicircum{}s/k) matching their E-restrictions are nonempty clopen with the finite-intersection property. Profinite compactness gives one gamma matching every coordinate. Thus full infinite fibres are orbits, without extending through an original fraction field across support.

Every Galois-stable ideal descends without division by group order. In a finite Galois E/k choose basis e\_j and trace-dual c\_i. Nondegeneracy follows from the invertible embedding matrix M: its trace Gram matrix is M\textasciicircum{}t M and determinant (det M)\textasciicircum{}2 is nonzero, including in characteristic dividing |G|. If s=sum r\_j e\_j is in a stable ideal, sum\_sigma sigma(c\_i s)=r\_i is in its contraction. This proves the ideal is the extension of its contraction. At the infinite stage choose a finite Galois field containing finitely many ideal generators; faithful flat contraction and extension of finite-field automorphisms make that finite ideal stable, so it descends too. Canonical ideals B\_m for a pulled obstacle are stable and contract to A\_m. Therefore B\_m=A\_m R\_s in EVERY degree, their profiles and raw envelopes pull back exactly, and uniform errors are unchanged. Continuity is equivalent by the compact quotient. FG equivalence follows by faithfully flat detection of degree-generation multiplication maps. For invariant C, approximate by a rational primary profile at a finite stage and replace its ideal by the sum of its Galois translates; its maximum profile stays within the same uniform error and descends by trace-dual extraction. No arbitrary invariant infinite D descent is claimed.

In characteristic zero k\textasciicircum{}s is algebraically closed. Finite linear smooth quotients at the rational origin over any characteristic-zero k satisfy the required geometric hypotheses after every field extension: polynomial invariants commute with extension degree by degree; orbit polynomials give a finite invariant coefficient algebra and a finite smooth cover; their nonleading coefficients vanish at the origin and force its fibre to have the unique origin; invariant fractions and integral closure prove normality after extension. The smooth quotient result applies only over k\textasciicircum{}s, and the every-degree identity descends it over k with unchanged model error constants. This includes k[a,b,c,z]/(ac-b\textasciicircum{}2) at the origin, char(k)=0, a=x\textasciicircum{}2,b=xy,c=y\textasciicircum{}2, with its normal nonisolated singular locus (a,b,c). It does not prove positive-characteristic regular source CE.

For an injective finite local excellent reduced characteristic-p extension R subset S with S\textasciicircum{}q subset R, q=p\textasciicircum{}e, every semivaluation has the unique lift w(s)=v(s\textasciicircum{}q)/q. Frobenius addition/multiplication proves the axioms, powers determine every support, and aS normalization proves centredness; the coordinate formulas prove a homeomorphism. For I=(s\_i) primary or unit, J=(s\_i\textasciicircum{}q) in R obeys v(J)=q w(I). If I contains (aS)\textasciicircum{}N, raise each actual expression x=sum s\_i u\_i, x in a\textasciicircum{}N, to obtain x\textasciicircum{}q in J; thus J is primary, and raising a unit expression gives the unit case. Rational density identifies C isometrically; actual decreasing nets identify D. Hence ALL pointwise raw envelopes commute, and CE is equivalent. Canonical contraction is A\_m=B\_m intersect R and J\_m subset A\_(qm), giving the conditional rate (2qE+qM\_1)/k for k>=2q by the same residue argument. Same-degree scalar extension is not asserted.

If S is a Noetherian local union of finite local root subextensions and its maximal ideal/gauge are extended, each element has some p-power in R. Common exponents prove the coordinate formula, and each finitely generated ideal plus finitely many primary-containment coefficients has one common exponent. The same primary-profile argument identifies C,D and envelopes, but supplies no global rate when exponents are unbounded. Finite purely inseparable extensions of normal domains satisfy S\textasciicircum{}q subset R by normality and a field exponent. Geometrically integral/normal rational finite-type germs tensor k\textasciicircum{}perf meet the union hypotheses by Frobenius on their finitely many tensor coefficients. Combining separable closure and perfection reduces their compatible obstacle problem to algebraic closure, without supplying characteristic-p CE there. Mixed characteristic has no such formula.

No literal mathematical correction to these current413 scoped proofs was needed. The real-parameter BFJ step and support-quotient trace are essential, and were checked rather than inferred from old verdicts. Its finite diagnostics and compilation remain observations. Historical recursive audit is not inherited from retained text. Artifacts: n2.txt, distinct expanded n2.tex, source-application-reading-ledger.md and proof-companion.md. Rerun: python build\_delivery.py; python check\_delivery.py.

\subsection*{Expanded scope-sensitive transport identities}
On the published smooth source, using one fixed compact complex,
\[
 W_t\le W\le(1-C/t)W_{t-C},\qquad t=m+C,
 \qquad (1+C/m)W\le W_m\le W.
\]
All-retraction infima give the all-support estimate
$0\le H-h_m\le CM/m$.
The separable coefficient identity carries the actual degree:
\[
 A_m^{R_{k^s}}(\rho^*f)=A_m^R(f)R_{k^s},\qquad
 h_m^{R_{k^s}}=\rho^*h_m^R,\qquad
 p_{R_{k^s}}(\rho^*f)=\rho^*p_R(f).
\]
The Frobenius-root formula instead gives
\[
 w(s)=q^{-1}v(s^q),\qquad v(J_m)=q\,w(B_m),\qquad
 J_m\subset A_{qm}.
\]
It does not assert the same-degree scalar-extension identity.
For a source bound $E/m$ and $k\ge2q$ the target bound is
$0\le p_R-h_k^R\le(2qE+qM_1)/k$.
